% generated by GAPDoc2LaTeX from XML source (Frank Luebeck)
\documentclass[a4paper,11pt]{report}

\usepackage[top=37mm,bottom=37mm,left=27mm,right=27mm]{geometry}
\sloppy
\pagestyle{myheadings}
\usepackage{amssymb}
\usepackage[utf8]{inputenc}
\usepackage{makeidx}
\makeindex
\usepackage{color}
\definecolor{FireBrick}{rgb}{0.5812,0.0074,0.0083}
\definecolor{RoyalBlue}{rgb}{0.0236,0.0894,0.6179}
\definecolor{RoyalGreen}{rgb}{0.0236,0.6179,0.0894}
\definecolor{RoyalRed}{rgb}{0.6179,0.0236,0.0894}
\definecolor{LightBlue}{rgb}{0.8544,0.9511,1.0000}
\definecolor{Black}{rgb}{0.0,0.0,0.0}

\definecolor{linkColor}{rgb}{0.0,0.0,0.554}
\definecolor{citeColor}{rgb}{0.0,0.0,0.554}
\definecolor{fileColor}{rgb}{0.0,0.0,0.554}
\definecolor{urlColor}{rgb}{0.0,0.0,0.554}
\definecolor{promptColor}{rgb}{0.0,0.0,0.589}
\definecolor{brkpromptColor}{rgb}{0.589,0.0,0.0}
\definecolor{gapinputColor}{rgb}{0.589,0.0,0.0}
\definecolor{gapoutputColor}{rgb}{0.0,0.0,0.0}

%%  for a long time these were red and blue by default,
%%  now black, but keep variables to overwrite
\definecolor{FuncColor}{rgb}{0.0,0.0,0.0}
%% strange name because of pdflatex bug:
\definecolor{Chapter }{rgb}{0.0,0.0,0.0}
\definecolor{DarkOlive}{rgb}{0.1047,0.2412,0.0064}


\usepackage{fancyvrb}

\usepackage{mathptmx,helvet}
\usepackage[T1]{fontenc}
\usepackage{textcomp}


\usepackage[
            pdftex=true,
            bookmarks=true,        
            a4paper=true,
            pdftitle={Written with GAPDoc},
            pdfcreator={LaTeX with hyperref package / GAPDoc},
            colorlinks=true,
            backref=page,
            breaklinks=true,
            linkcolor=linkColor,
            citecolor=citeColor,
            filecolor=fileColor,
            urlcolor=urlColor,
            pdfpagemode={UseNone}, 
           ]{hyperref}

\newcommand{\maintitlesize}{\fontsize{50}{55}\selectfont}

% write page numbers to a .pnr log file for online help
\newwrite\pagenrlog
\immediate\openout\pagenrlog =\jobname.pnr
\immediate\write\pagenrlog{PAGENRS := [}
\newcommand{\logpage}[1]{\protect\write\pagenrlog{#1, \thepage,}}
%% were never documented, give conflicts with some additional packages

\newcommand{\GAP}{\textsf{GAP}}

%% nicer description environments, allows long labels
\usepackage{enumitem}
\setdescription{style=nextline}

%% depth of toc
\setcounter{tocdepth}{1}





%% command for ColorPrompt style examples
\newcommand{\gapprompt}[1]{\color{promptColor}{\bfseries #1}}
\newcommand{\gapbrkprompt}[1]{\color{brkpromptColor}{\bfseries #1}}
\newcommand{\gapinput}[1]{\color{gapinputColor}{#1}}


\begin{document}

\logpage{[ 0, 0, 0 ]}
\begin{titlepage}
\mbox{}\vfill

\begin{center}{\maintitlesize \textbf{ FPLSA \mbox{}}}\\
\vfill

\hypersetup{pdftitle= FPLSA }
\markright{\scriptsize \mbox{}\hfill  FPLSA  \hfill\mbox{}}
{\Huge \textbf{ Finitely Presented Lie Algebras \mbox{}}}\\
\vfill

{\Huge  1.2.8 \mbox{}}\\[1cm]
{ 26 December 2025 \mbox{}}\\[1cm]
\mbox{}\\[2cm]
{\Large \textbf{ Vladimir Gerdt\\
   \mbox{}}}\\
{\Large \textbf{ Vladimir Kornyak\\
   \mbox{}}}\\
\hypersetup{pdfauthor= Vladimir Gerdt\\
   ;  Vladimir Kornyak\\
   }
\end{center}\vfill

\mbox{}\\
{\mbox{}\\
\small \noindent \textbf{ Vladimir Gerdt\\
   }  Email: \href{mailto://gerdt@jinr.ru} {\texttt{gerdt@jinr.ru}}\\
  Homepage: \href{http://compalg.jinr.ru/CAGroup/Gerdt/} {\texttt{http://compalg.jinr.ru/CAGroup/Gerdt/}}}\\
{\mbox{}\\
\small \noindent \textbf{ Vladimir Kornyak\\
   }  Email: \href{mailto://vkornyak@gmail.com} {\texttt{vkornyak@gmail.com}}\\
  Homepage: \href{http://compalg.jinr.ru/CAGroup/Kornyak} {\texttt{http://compalg.jinr.ru/CAGroup/Kornyak}}}\\
\end{titlepage}

\newpage\setcounter{page}{2}
\newpage

\def\contentsname{Contents\logpage{[ 0, 0, 1 ]}}

\tableofcontents
\newpage

 
\chapter{\textcolor{Chapter }{The FPLSA Package}}\label{The FPLSA Package}
\logpage{[ 1, 0, 0 ]}
\hyperdef{L}{X805DE95983EEA672}{}
{
  This chapter describes the \textsf{GAP} package \textsf{FPLSA} , an interface to the \texttt{fplsa} program by V.{\nobreakspace}Gerdt and V.{\nobreakspace}Kornyak (version 4) for
the computation with finitely presented Lie superalgebras. 

 At present \textsf{GAP} uses only the facility to compute a structure constants table of a finite
dimensional Lie algebra over the rationals that is given by a finite
presentation. 

 The package uses an external binary, probably it will only work on UNIX
platforms. 

  
\section{\textcolor{Chapter }{Main Functions}}\label{Main Functions}
\logpage{[ 1, 1, 0 ]}
\hyperdef{L}{X8240F18D7AFA30F4}{}
{
  A finitely presented Lie algebra is a quotient of a free Lie algebra by an
ideal generated by a finite number of elements. In \textsf{GAP} a free Lie algebra can be created by the command \texttt{FreeLieAlgebra}; we refer to the \textsf{GAP} Reference Manual for more details. A finitely presented Lie algebra \mbox{\texttt{\mdseries\slshape K}} can be constructed by \texttt{\mbox{\texttt{\mdseries\slshape K}} := \mbox{\texttt{\mdseries\slshape L}}/\mbox{\texttt{\mdseries\slshape rels}}}, where \mbox{\texttt{\mdseries\slshape L}} is a free Lie algebra and \mbox{\texttt{\mdseries\slshape rels}} a list of elements of \mbox{\texttt{\mdseries\slshape L}} that constitute the relations that hold in \mbox{\texttt{\mdseries\slshape K}}. Given a finitely presented Lie algebra we want to calculate a basis and a
multiplication table of it. The interface to the \textsf{FPLSA} package comes with two related functions for doing that. 

 

\subsection{\textcolor{Chapter }{SCAlgebraInfoOfFpLieAlgebra}}
\logpage{[ 1, 1, 1 ]}\nobreak
\hyperdef{L}{X8019B72B8298566A}{}
{\noindent\textcolor{FuncColor}{$\triangleright$\enspace\texttt{SCAlgebraInfoOfFpLieAlgebra({\mdseries\slshape L, rels, limit, words, rels})\index{SCAlgebraInfoOfFpLieAlgebra@\texttt{SCAlgebraInfoOfFpLieAlgebra}}
\label{SCAlgebraInfoOfFpLieAlgebra}
}\hfill{\scriptsize (function)}}\\


 Let \mbox{\texttt{\mdseries\slshape L}} be a free Lie algebra over the rationals, \mbox{\texttt{\mdseries\slshape rels}} a list of elements in \mbox{\texttt{\mdseries\slshape L}}, \mbox{\texttt{\mdseries\slshape limit}} a positive integer and \mbox{\texttt{\mdseries\slshape words}} and \mbox{\texttt{\mdseries\slshape rels}} two booleans. 

 If the algebra \texttt{\mbox{\texttt{\mdseries\slshape L}} / \mbox{\texttt{\mdseries\slshape rels}}} is finite\texttt{\symbol{45}}dimensional and if a basis of this algebra can be
constructed using elements in \mbox{\texttt{\mdseries\slshape L}} that involve only words of length at most \mbox{\texttt{\mdseries\slshape limit}} then \texttt{SCAlgebraInfoOfFpLieAlgebra} returns a record whose component \texttt{sc} contains an algebra that is isomorphic with \texttt{\mbox{\texttt{\mdseries\slshape L}} / \mbox{\texttt{\mdseries\slshape rels}}}. Otherwise \texttt{fail} is returned. 

 The function calls the \texttt{fplsa} standalone. 

 If \mbox{\texttt{\mdseries\slshape words}} is \texttt{true} then the component \texttt{words} of the result record contains a list of elements in \mbox{\texttt{\mdseries\slshape L}} that correspond to the basis elements. 

 If \mbox{\texttt{\mdseries\slshape rels}} is \texttt{true} then the component \texttt{rels} of the result record contains a list of reduced relators in \mbox{\texttt{\mdseries\slshape L}} that describes how algebra generators of \mbox{\texttt{\mdseries\slshape L}} are expressed in terms of the basis elements. 

 
\begin{Verbatim}[commandchars=!@|,fontsize=\small,frame=single,label=Example]
  !gapprompt@gap>| !gapinput@LoadPackage( "fplsa" );|
  true
  !gapprompt@gap>| !gapinput@L:= FreeLieAlgebra( Rationals, "x", "y" );|
  <Lie algebra over Rationals, with 2 generators>
  !gapprompt@gap>| !gapinput@g:= GeneratorsOfAlgebra( L );; x:= g[1];; y:= g[2];;|
  !gapprompt@gap>| !gapinput@rels:= [ x*(x*y) - x*y, y*(y*(x*y)) ];|
  [ (-1)*(x*y)+(1)*(x*(x*y)), (1)*(((x*y)*y)*y) ]
  !gapprompt@gap>| !gapinput@SCAlgebraInfoOfFpLieAlgebra( L, rels, 100, true, true );|
  rec( sc := <Lie algebra of dimension 4 over Rationals>, 
    words := [ (1)*x, (1)*y, (1)*(y*x), (1)*((y*x)*y) ], 
    rels := [ (1)*((y*x)*x)+(1)*(y*x), (1)*(((y*x)*y)*y), (1)*(((y*x)*y)*(y*x)) 
       ] )
\end{Verbatim}
 }

 

\subsection{\textcolor{Chapter }{ IsomorphicSCAlgebra}}
\logpage{[ 1, 1, 2 ]}\nobreak
\hyperdef{L}{X7C5E1F7B80268EBD}{}
{\noindent\textcolor{FuncColor}{$\triangleright$\enspace\texttt{ IsomorphicSCAlgebra({\mdseries\slshape K[, bound]})\index{ IsomorphicSCAlgebra@\texttt{ IsomorphicSCAlgebra}}
\label{ IsomorphicSCAlgebra}
}\hfill{\scriptsize (function)}}\\


 computes a Lie algebra given by a multiplication table isomorphic to the
finitely presented Lie algebra \mbox{\texttt{\mdseries\slshape K}}. If the optional parameter \mbox{\texttt{\mdseries\slshape bound}} is specified the computation will be carried out using monomials of degree at
most \mbox{\texttt{\mdseries\slshape bound}}. If \mbox{\texttt{\mdseries\slshape bound}} is not specified, then it will initially be set to 10000. If this does not
suffice to calculate a multiplication table of the algebra, then the bound
will be increased until a multiplication table is found. If the computation
was successful then a structure constants algebra will be returned isomorphic
to \mbox{\texttt{\mdseries\slshape K}}. Otherwise \texttt{fail} will be returned. 

 
\begin{Verbatim}[commandchars=!@|,fontsize=\small,frame=single,label=Example]
  !gapprompt@gap>| !gapinput@LoadPackage( "fplsa" );|
  true
  !gapprompt@gap>| !gapinput@L:= FreeLieAlgebra( Rationals, "x", "y" );|
  <Lie algebra over Rationals, with 2 generators>
  !gapprompt@gap>| !gapinput@g:= GeneratorsOfAlgebra( L );; x:= g[1];; y:= g[2];;|
  !gapprompt@gap>| !gapinput@rels:= [ x*(x*y) - x*y, y*(y*(x*y)) ];|
  [ (-1)*(x*y)+(1)*(x*(x*y)), (1)*(((x*y)*y)*y) ]
  !gapprompt@gap>| !gapinput@K:= L/rels;|
  <Lie algebra over Rationals, with 2 generators>
  !gapprompt@gap>| !gapinput@IsomorphicSCAlgebra( K );|
  <Lie algebra of dimension 4 over Rationals>
\end{Verbatim}
 }

 }

  
\section{\textcolor{Chapter }{Auxiliary Variables of FPLSA}}\label{Auxiliary Variables of FPLSA}
\logpage{[ 1, 2, 0 ]}
\hyperdef{L}{X7C991B3B79A45BE5}{}
{
  

\subsection{\textcolor{Chapter }{FPLSA}}
\logpage{[ 1, 2, 1 ]}\nobreak
\hyperdef{L}{X87E75AD286BA918B}{}
{\noindent\textcolor{FuncColor}{$\triangleright$\enspace\texttt{FPLSA\index{FPLSA@\texttt{FPLSA}}
\label{FPLSA}
}\hfill{\scriptsize (global variable)}}\\


 

 is the global record used by the functions in the package. Besides components
that describe parameters for the standalone, the following components are
used. 
\begin{description}
\item[{\texttt{Relation{\textunderscore}size} }]  parameter that controls the memory usage by the fplsa program, 
\item[{\texttt{Lie{\textunderscore}monomial{\textunderscore}table{\textunderscore}size} }]  parameter that controls the memory usage by the fplsa program, 
\item[{\texttt{Node{\textunderscore}Lie{\textunderscore}term{\textunderscore}size} }]  parameter that controls the memory usage by the fplsa program, 
\item[{\texttt{Node{\textunderscore}scalar{\textunderscore}factor{\textunderscore}size} }]  parameter that controls the memory usage by the fplsa program, 
\item[{\texttt{Node{\textunderscore}scalar{\textunderscore}term{\textunderscore}size} }]  parameter that controls the memory usage by the fplsa program, 
\item[{\texttt{progname} }]  the file name of the executable, 
\item[{\texttt{T} }]  structure constants table of the algebra under consideration, 
\item[{\texttt{words} }]  list of elements in the free Lie algebra that correspond to the basis
elements, 
\item[{\texttt{rels} }]  list of relators in the free Lie algebra that are used to express redundant
algebra generators in terms of the basis. 
\end{description}
 

 In order to be able to run the \texttt{fplsa} program successfully, it might be necessary to customize the parameters that
control the memory the the program uses, to suit the computer the user is
working on. In particular if the program exits with an error message of the
form: \texttt{Error, the process did not succeed}, then it may be necessary to change these parameters. 

 
\begin{Verbatim}[commandchars=!@|,fontsize=\small,frame=single,label=Example]
  !gapprompt@gap>| !gapinput@LoadPackage( "fplsa" );|
  true
  !gapprompt@gap>| !gapinput@L:= FreeLieAlgebra( Rationals, "x", "y" );;|
  !gapprompt@gap>| !gapinput@g:= GeneratorsOfAlgebra( L );; x:= g[1];; y:= g[2];;|
  !gapprompt@gap>| !gapinput@rels:= [ x*(x*y) - x*y, y*(y*(x*y)) ];;|
  !gapprompt@gap>| !gapinput@SCAlgebraInfoOfFpLieAlgebra( L, rels, 100, true, true );;|
  !gapprompt@gap>| !gapinput@FPLSA;|
  rec( Relation_size := 2500000, Lie_monomial_table_size := 1000000, 
    Node_Lie_term_size := 2000000, Node_scalar_factor_size := 2000, 
    Node_scalar_term_size := 20000, progname := "fplsa4", 
    T := [ [ [ [  ], [  ] ], [ [ 3 ], [ -1 ] ], [ [ 3 ], [ 1 ] ], 
            [ [ 4 ], [ 1 ] ] ], 
        [ [ [ 3 ], [ 1 ] ], [ [  ], [  ] ], [ [ 4 ], [ -1 ] ], [ [  ], [  ] ] ],
        [ [ [ 3 ], [ -1 ] ], [ [ 4 ], [ 1 ] ], [ [  ], [  ] ], [ [  ], [  ] ] ],
        [ [ [ 4 ], [ -1 ] ], [ [  ], [  ] ], [ [  ], [  ] ], [ [  ], [  ] ] ], 
        -1, 0 ], words := [ 1, 2, [ 2, 1 ], [ [ 2, 1 ], 2 ] ], 
    rels := [ [ [ [ 2, 1 ], 1 ], 1, [ 2, 1 ], 1 ], 
        [ [ [ [ 2, 1 ], 2 ], 2 ], 1 ], [ [ [ [ 2, 1 ], 2 ], [ 2, 1 ] ], 1 ] ] )
\end{Verbatim}
 }

 }

  
\section{\textcolor{Chapter }{Installing the FPLSA Package}}\label{Installing the FPLSA Package}
\logpage{[ 1, 3, 0 ]}
\hyperdef{L}{X7F15E49B8173004C}{}
{
  To install unpack the archive file in a directory in the \texttt{pkg} hierarchy of your version of \textsf{GAP}. (This might be the \texttt{pkg} directory of the \textsf{GAP} home directory; it is however also possible to keep an additional \texttt{pkg} directory in your private directories, see Section{\nobreakspace} (\textbf{Reference: Installing a GAP Package}) of the \textsf{GAP} Reference Manual for details on how to do this.) Go to the newly created \texttt{fplsa} directory and call \texttt{./configure \mbox{\texttt{\mdseries\slshape path}}} where \mbox{\texttt{\mdseries\slshape path}} is the path to the \textsf{GAP} home directory. So for example if you install the package in the main \texttt{pkg} directory call 
\begin{Verbatim}[commandchars=!@|,fontsize=\small,frame=single,label=Example]
  ./configure ../..
\end{Verbatim}
 This will fetch the architecture type for which \textsf{GAP} has been compiled last and create a \texttt{Makefile}. Now simply call 
\begin{Verbatim}[commandchars=!@|,fontsize=\small,frame=single,label=Example]
  make
\end{Verbatim}
 to compile the binary and to install it in the appropriate place. 

 If you use this installation of \textsf{GAP} on different hardware platforms you will have to compile the binary for each
platform separately. This is done by calling \texttt{configure} and \texttt{make} for the package anew immediately after compiling \textsf{GAP} itself for the respective architecture. If your version of \textsf{GAP} is already compiled (and has last been compiled on the same architecture) you
do not need to compile \textsf{GAP} again, it is sufficient to call the \texttt{configure} script in the \textsf{GAP} home directory. 

    

 }

 }

 \def\indexname{Index\logpage{[ "Ind", 0, 0 ]}
\hyperdef{L}{X83A0356F839C696F}{}
}

\cleardoublepage
\phantomsection
\addcontentsline{toc}{chapter}{Index}


\printindex

\immediate\write\pagenrlog{["Ind", 0, 0], \arabic{page},}
\newpage
\immediate\write\pagenrlog{["End"], \arabic{page}];}
\immediate\closeout\pagenrlog
\end{document}
